
The main result of this article has been to show that the existential
theory of generalised Sem{\"e}nov arithmetic is decidable in
EXPSPACE. As an application of this result, we showed that a highly
relevant class of string constraints with length constraints is also
decidable in EXPSPACE; the decidability of this class was the main
problem left open in~\cite{BDGKMMN23}. On a technical level, those
results were obtained by showing that a restricted class of labelled
affine VASS has an EXPSPACE-decidable language emptiness problem. The
structural restrictions imposed on those restricted \laVASS are rather
strong, though necessary to obtain a decidable class of \laVASS.

An interesting aspect of our approach is that it establishes
automaticity of the existential fragment of a logical theory that is
different from traditional notions of automaticity, which are based on
finite-state automata or tree automata over finite or infinite words
and trees~\cite{KN95,BG00}, respectively. It would be interesting to
better understand whether there are natural logical theories whose
(existential) fragments are, say, Petri-net or visibly-pushdown
automatic.

We have ignored algorithmic lower bounds throughout this article, but
it would, of course, be interesting to see whether the upper bounds of
the decision problems we considered in this article are tight. It is
clear that generalised Sem{\"e}nov arithmetic is PSPACE-hard since it
can readily express the DFA intersection non-emptiness problem, but
this still leaves a considerable gap with respect to the EXPSPACE
upper bound we established. In particular, the recent results
of~\cite{BCM23} showing an NEXP upper bound for the existential
fragment of Sem{\"e}nov arithmetic suggest that, if an EXPSPACE lower
bound for existential generalised Sem{\"e}nov arithmetic is possible,
it will require the use of regular predicates.
